\clearpage
\chapter*{Erratum to Exposé no. 182}
\addcontentsline{toc}{chapter}{Erratum to Exposé no. 182}
\setcounter{footnote}{1}
\src{298}

\noindent\textbf{2. Formal geometry and algebraic geometry.}

\begin{center}
[Bourbaki Seminar, vol. 11, 1958/59, no. 182, 28 pp.]
\end{center}

\noindent\textbf{Page 8, Theorem 6.} --- The hypothesis “$f$ is
quasi-projective” may be replaced by the weaker hypothesis “$f$ is separated,”
by virtue of the following result (cf. SGA VIII, 6.2): every quasi-finite
separated morphism $f:X\longrightarrow Y$ is quasi-projective.

\src{299}
\noindent\textbf{Page 12, and page 14, Remark 1.} --- J.-P. Serre constructed
a nonsingular projective variety of dimension 3 over an algebraically closed
field $k$ of characteristic $p>0$ which is not obtained by reduction from a
proper scheme over an integral local ring with residue field $k$ and fraction
field of characteristic zero.\footnote{Serre (Jean-Pierre). --- Exemples de
variétés projectives en caractéristique $p$ non relevables en caractéristique
zéro, \emph{Proc. Nat. Acad. Sc. U.S.A.}, vol. 47, 1961, pp. 108--109.} Mumford
is said to have found an analogous example involving a nonsingular projective
surface.

\noindent\textbf{Page 15, Remarks 2 and 3.} --- I am now less optimistic about
the results conjectured here. Nevertheless, the problem concerning variations
of structure for projective space, mentioned at the end of Remark 3, has been
solved affirmatively by Hironaka, and the analogous problem for abelian
varieties has been solved by Koizumi.

\noindent\textbf{Page 16, line 3.} --- Mumford recently constructed the moduli
schemes for curves of genus $g$ (cf. SHMI). Theorem 10 also proves that the
“Jacobi schemes of level $n$” in the theory of moduli are nonsingular (and in
fact smooth over $\mathbf Z$).

\noindent\textbf{Page 23, line 12, first line of Corollary 5.} --- Add: “$X$
or $Y$ being proper over $k$.”

The substance of §§ 1 through 5 is contained in the published part of EGA III,
and that of §§ 6 and 7 is contained in SGA III. For the study of the
fundamental group, see SGA V, IX, X, XI, as well as SGA 1962 (Exposés X, XII,
and XIII) for Lefschetz-type theorems and numerous open questions. Only the
theory of tamely ramified coverings (cf. Theorem 14) has not yet received a
formal exposition. The corollary to Theorem 14, which gives the complete
determination of Galois coverings of order prime to the characteristic of an
algebraic curve over an algebraically closed field, has been used essentially
on three occasions:

1\textsuperscript{o} in Igusa's proof of the Picard inequality for
nonsingular projective surfaces in arbitrary characteristic;

2\textsuperscript{o} in the study (developed independently by Ogg and
Šafarevič) of the group of homogeneous principal bundles under an abelian
variety defined over a one-variable function field, in arbitrary
characteristic;

3\textsuperscript{o} in Artin's recent proof of certain key theorems on the
“Weil cohomology” of algebraic varieties.
